\documentclass[12pt,a4paper]{article}
\usepackage[utf8]{inputenc}
\usepackage{amsmath, amssymb}
\usepackage{enumitem}
\usepackage[margin=2cm]{geometry}

\begin{document}

\begin{center}
    \textbf{\Large Latihan Soal Matematika SMK Kelas X} \\
    \textbf{Persiapan Ujian}
\end{center}
\vspace{0.5cm}

\begin{enumerate}[leftmargin=*]

% 1. Eksponen (menyatakan/menyederhanakan)
\item Bentuk sederhana dari $\frac{a^4 \cdot b^3}{a^2 \cdot b}$ adalah \dots
    \begin{enumerate}[label=\alph*.]
        \item $a^6 b^4$
        \item $a^2 b^2$
        \item $a^2 b$
        \item $a^8 b^3$
        \item $a b^2$
    \end{enumerate}

\item Nilai dari $(2^3)^2 \cdot 2^{-4}$ adalah \dots
    \begin{enumerate}[label=\alph*.]
        \item $2$
        \item $4$
        \item $8$
        \item $16$
        \item $32$
    \end{enumerate}

\item Bentuk pangkat positif dari $3x^{-2}y^3$ adalah \dots
    \begin{enumerate}[label=\alph*.]
        \item $\frac{3y^3}{x^2}$
        \item $\frac{y^3}{3x^2}$
        \item $\frac{3}{x^2 y^3}$
        \item $\frac{x^2}{3y^3}$
        \item $3x^2 y^3$
    \end{enumerate}

% 2. Menentukan nilai logaritma dasar
\item Nilai dari $^2\log 16$ adalah \dots
    \begin{enumerate}[label=\alph*.]
        \item $2$
        \item $3$
        \item $4$
        \item $8$
        \item $16$
    \end{enumerate}

\item Jika $^3\log 27 = x$, maka nilai $x$ adalah \dots
    \begin{enumerate}[label=\alph*.]
        \item $2$
        \item $3$
        \item $4$
        \item $9$
        \item $27$
    \end{enumerate}

\item Hasil dari $\log 100 + \log 10$ adalah \dots
    \begin{enumerate}[label=\alph*.]
        \item $1$
        \item $2$
        \item $3$
        \item $10$
        \item $110$
    \end{enumerate}

% 3. Menyelesaikan spldv
\item Himpunan penyelesaian dari persamaan $x + y = 5$ dan $x - y = 1$ adalah \dots
    \begin{enumerate}[label=\alph*.]
        \item $\{(2, 3)\}$
        \item $\{(3, 2)\}$
        \item $\{(4, 1)\}$
        \item $\{(1, 4)\}$
        \item $\{(5, 0)\}$
    \end{enumerate}

\item Jika $2x + y = 7$ dan $x - y = 2$, maka nilai $x + y$ adalah \dots
    \begin{enumerate}[label=\alph*.]
        \item $3$
        \item $4$
        \item $5$
        \item $6$
        \item $7$
    \end{enumerate}

% 4. Menyelesaikan persamaan linear 1 variabel
\item Nilai $x$ yang memenuhi persamaan $3x - 5 = 10$ adalah \dots
    \begin{enumerate}[label=\alph*.]
        \item $3$
        \item $4$
        \item $5$
        \item $10$
        \item $15$
    \end{enumerate}

\item Penyelesaian dari persamaan $2(x - 3) = x + 4$ adalah \dots
    \begin{enumerate}[label=\alph*.]
        \item $x = -2$
        \item $x = 2$
        \item $x = 5$
        \item $x = 10$
        \item $x = 12$
    \end{enumerate}

% 5. Menyelesaikan spltv
\item Diketahui sistem persamaan $x+y+z=6$, $2x-y+z=3$, dan $x+2y-z=2$. Nilai $x$ adalah \dots
    \begin{enumerate}[label=\alph*.]
        \item $1$
        \item $2$
        \item $3$
        \item $4$
        \item $5$
    \end{enumerate}

\item Pada SPLTV dengan variabel $a, b, c$, jika diketahui $a=2, b=1$, dan $a+b+c=7$, maka nilai $c$ adalah \dots
    \begin{enumerate}[label=\alph*.]
        \item $2$
        \item $3$
        \item $4$
        \item $5$
        \item $6$
    \end{enumerate}

% 6. Bentuk umum fungsi kuadrat
\item Bentuk umum dari fungsi kuadrat adalah \dots
    \begin{enumerate}[label=\alph*.]
        \item $f(x) = ax + b$
        \item $f(x) = ax^2 + bx + c$
        \item $f(x) = ax^3 + bx^2 + cx + d$
        \item $f(x) = \frac{a}{x} + b$
        \item $f(x) = a^x + b$
    \end{enumerate}

\item Pada fungsi kuadrat $f(x) = -2x^2 + 3x - 5$, nilai $a, b,$ dan $c$ berturut-turut adalah \dots
    \begin{enumerate}[label=\alph*.]
        \item $2, 3, 5$
        \item $-2, 3, 5$
        \item $-2, 3, -5$
        \item $2, -3, -5$
        \item $-2, -3, -5$
    \end{enumerate}

% 7. Menentukan sumbu simetri
\item Persamaan sumbu simetri dari grafik fungsi kuadrat $f(x) = x^2 - 4x + 3$ adalah \dots
    \begin{enumerate}[label=\alph*.]
        \item $x = -4$
        \item $x = -2$
        \item $x = 2$
        \item $x = 3$
        \item $x = 4$
    \end{enumerate}

\item Sumbu simetri dari fungsi $f(x) = -2x^2 + 8x - 1$ adalah \dots
    \begin{enumerate}[label=\alph*.]
        \item $x = 2$
        \item $x = -2$
        \item $x = 4$
        \item $x = -4$
        \item $x = 8$
    \end{enumerate}

% 8. Menganalisis titik puncak dan arah parabola
\item Grafik fungsi kuadrat $f(x) = -x^2 + 2x + 8$ akan terbuka ke \dots dan memiliki titik \dots
    \begin{enumerate}[label=\alph*.]
        \item atas, minimum
        \item atas, maksimum
        \item bawah, minimum
        \item bawah, maksimum
        \item samping, minimum
    \end{enumerate}

\item Nilai minimum dari fungsi kuadrat $f(x) = x^2 - 6x + 10$ adalah \dots
    \begin{enumerate}[label=\alph*.]
        \item $1$
        \item $2$
        \item $3$
        \item $4$
        \item $10$
    \end{enumerate}

% 9. Menentukan langkah langkah menggambar grafik fungsi kuadrat
\item Langkah pertama yang paling tepat dalam menggambar grafik fungsi kuadrat adalah \dots
    \begin{enumerate}[label=\alph*.]
        \item Menentukan titik potong dengan sumbu x
        \item Menentukan titik puncak parabola
        \item Menarik kurva mulus
        \item Menentukan persamaan sumbu simetri
        \item Menentukan titik potong dengan sumbu y
    \end{enumerate}

\item Titik potong grafik fungsi $f(x) = x^2 - x - 6$ dengan sumbu Y adalah \dots
    \begin{enumerate}[label=\alph*.]
        \item $(0, 6)$
        \item $(0, -6)$
        \item $(6, 0)$
        \item $(-6, 0)$
        \item $(0, 1)$
    \end{enumerate}

% 10. Menentukan nilai sin, cos, tan (sudut istimewanya)
\item Nilai dari $\sin 30^\circ$ adalah \dots
    \begin{enumerate}[label=\alph*.]
        \item $0$
        \item $\frac{1}{2}$
        \item $\frac{1}{2}\sqrt{2}$
        \item $\frac{1}{2}\sqrt{3}$
        \item $1$
    \end{enumerate}

\item Nilai dari $\cos 60^\circ + \sin 30^\circ$ adalah \dots
    \begin{enumerate}[label=\alph*.]
        \item $0$
        \item $\frac{1}{2}$
        \item $1$
        \item $\sqrt{2}$
        \item $\sqrt{3}$
    \end{enumerate}

\item Nilai dari $\tan 45^\circ$ adalah \dots
    \begin{enumerate}[label=\alph*.]
        \item $0$
        \item $\frac{1}{3}\sqrt{3}$
        \item $1$
        \item $\sqrt{3}$
        \item Tidak terdefinisi
    \end{enumerate}

% 11. Menghitung panjang sisi dengan perbandingan trigonometri
\item Pada segitiga siku-siku, jika panjang hipotenusa 10 cm dan salah satu sudutnya $30^\circ$, panjang sisi di depan sudut tersebut adalah \dots
    \begin{enumerate}[label=\alph*.]
        \item $5$ cm
        \item $5\sqrt{2}$ cm
        \item $5\sqrt{3}$ cm
        \item $10$ cm
        \item $10\sqrt{3}$ cm
    \end{enumerate}

\item Diketahui segitiga siku-siku ABC siku-siku di B. Jika panjang AB = 3 cm dan BC = 4 cm, maka nilai $\sin A$ adalah \dots
    \begin{enumerate}[label=\alph*.]
        \item $\frac{3}{5}$
        \item $\frac{4}{5}$
        \item $\frac{3}{4}$
        \item $\frac{4}{3}$
        \item $\frac{5}{4}$
    \end{enumerate}

% 12. Sudut elevasi
\item Seseorang melihat puncak menara dengan sudut elevasi $45^\circ$. Jika jarak orang tersebut ke dasar menara adalah 20 m dan tinggi orang diabaikan, tinggi menara adalah \dots
    \begin{enumerate}[label=\alph*.]
        \item $10$ m
        \item $20$ m
        \item $20\sqrt{2}$ m
        \item $20\sqrt{3}$ m
        \item $40$ m
    \end{enumerate}

\item Sudut elevasi adalah sudut yang dibentuk oleh garis horizontal dan garis pandang pengamat ke arah \dots
    \begin{enumerate}[label=\alph*.]
        \item Bawah
        \item Samping kanan
        \item Samping kiri
        \item Atas
        \item Belakang
    \end{enumerate}

% 13. Sudut depresi
\item Pengamat di atas gedung setinggi 50 m melihat mobil di jalan dengan sudut depresi $30^\circ$. Jarak mobil ke dasar gedung adalah \dots
    \begin{enumerate}[label=\alph*.]
        \item $25$ m
        \item $50$ m
        \item $50\sqrt{3}$ m
        \item $100$ m
        \item $100\sqrt{3}$ m
    \end{enumerate}

\item Sudut depresi diukur dari garis horizontal ke arah \dots
    \begin{enumerate}[label=\alph*.]
        \item Atas objek
        \item Bawah (objek di bawah mata)
        \item Sejajar mata
        \item Titik pusat bumi
        \item Matahari
    \end{enumerate}

% 14. Menyajikan data (statistika)
\item Diagram yang paling tepat untuk menyajikan data persentase komponen dari suatu keseluruhan adalah \dots
    \begin{enumerate}[label=\alph*.]
        \item Diagram batang
        \item Diagram garis
        \item Diagram lingkaran
        \item Histogram
        \item Poligon frekuensi
    \end{enumerate}

\item Penyajian data distribusi frekuensi kumulatif dalam bentuk grafik disebut \dots
    \begin{enumerate}[label=\alph*.]
        \item Ogive
        \item Histogram
        \item Piktogram
        \item Poligon
        \item Kartogram
    \end{enumerate}

% 15. Menentukan panjang interval kelas
\item Jika jangkauan data adalah 40 dan banyak kelas yang diinginkan adalah 8, maka panjang interval kelas adalah \dots
    \begin{enumerate}[label=\alph*.]
        \item 4
        \item 5
        \item 6
        \item 7
        \item 8
    \end{enumerate}

\item Rumus Sturges untuk menentukan banyak kelas ($k$) dari $n$ datum adalah \dots
    \begin{enumerate}[label=\alph*.]
        \item $k = 1 + 3,3 \log n$
        \item $k = 1 - 3,3 \log n$
        \item $k = 3,3 + \log n$
        \item $k = 1 + \log 3,3n$
        \item $k = 3,3 \log (1+n)$
    \end{enumerate}

% 16. Menganalisis histogram
\item Pada histogram, sumbu mendatar (horizontal) biasanya menyatakan \dots
    \begin{enumerate}[label=\alph*.]
        \item Frekuensi
        \item Tepi kelas
        \item Nilai rata-rata
        \item Titik tengah kelas
        \item Frekuensi kumulatif
    \end{enumerate}

\item Luas persegi panjang pada histogram berbanding lurus dengan \dots kelas tersebut.
    \begin{enumerate}[label=\alph*.]
        \item Nilai rata-rata
        \item Frekuensi
        \item Batas bawah
        \item Jangkauan
        \item Simpangan
    \end{enumerate}

% 17. Mencari nilai modus data tunggal
\item Modus dari data: 7, 5, 8, 6, 7, 8, 7, 9, 5 adalah \dots
    \begin{enumerate}[label=\alph*.]
        \item 5
        \item 6
        \item 7
        \item 8
        \item 9
    \end{enumerate}

\item Diketahui data: 2, 3, 4, 3, 2, 5, 4, 3, 4, 4. Modus dari data tersebut adalah \dots
    \begin{enumerate}[label=\alph*.]
        \item 2
        \item 3
        \item 4
        \item 5
        \item 3 dan 4
    \end{enumerate}

% 18. Mencari rata rata data kelompok
\item Rumus untuk mencari rata-rata data berkelompok adalah \dots
    \begin{enumerate}[label=\alph*.]
        \item $\frac{\sum f_i}{\sum x_i}$
        \item $\frac{\sum f_i x_i}{\sum f_i}$
        \item $\sum (f_i + x_i)$
        \item $\frac{\sum x_i}{n}$
        \item $\sum f_i x_i$
    \end{enumerate}

\item Jika jumlah hasil kali frekuensi dengan titik tengah kelas adalah 1200 dan total frekuensi adalah 40, maka rata-ratanya adalah \dots
    \begin{enumerate}[label=\alph*.]
        \item 20
        \item 25
        \item 30
        \item 35
        \item 40
    \end{enumerate}

% 19. Mencari nilai modus data kelompok
\item Dalam rumus modus data berkelompok $Mo = Tb + \left(\frac{d_1}{d_1 + d_2}\right) p$, $Tb$ menyatakan \dots
    \begin{enumerate}[label=\alph*.]
        \item Tepi atas kelas modus
        \item Titik tengah kelas modus
        \item Batas bawah kelas modus
        \item Tepi bawah kelas modus
        \item Panjang interval
    \end{enumerate}

\item Jika tepi bawah kelas modus 49,5; panjang kelas 5; selisih frekuensi kelas modus dengan sebelumnya 4; dan dengan sesudahnya 6. Modusnya adalah \dots
    \begin{enumerate}[label=\alph*.]
        \item 50,5
        \item 51,0
        \item 51,5
        \item 52,0
        \item 52,5
    \end{enumerate}

% 20. Faktorial
\item Nilai dari $5!$ adalah \dots
    \begin{enumerate}[label=\alph*.]
        \item 15
        \item 24
        \item 60
        \item 120
        \item 720
    \end{enumerate}

\item Hasil dari $\frac{6!}{4!}$ adalah \dots
    \begin{enumerate}[label=\alph*.]
        \item 2
        \item 6
        \item 12
        \item 24
        \item 30
    \end{enumerate}

% 21. Permutasi
\item Banyaknya susunan huruf yang dapat dibentuk dari kata "BISA" adalah \dots
    \begin{enumerate}[label=\alph*.]
        \item 4
        \item 8
        \item 12
        \item 24
        \item 30
    \end{enumerate}

\item Nilai dari permutasi $P(5, 2)$ adalah \dots
    \begin{enumerate}[label=\alph*.]
        \item 10
        \item 15
        \item 20
        \item 25
        \item 60
    \end{enumerate}

\item Dari 6 orang siswa, akan dipilih ketua dan wakil ketua. Banyaknya cara pemilihannya adalah \dots
    \begin{enumerate}[label=\alph*.]
        \item 15
        \item 20
        \item 30
        \item 60
        \item 120
    \end{enumerate}

% 22. Menentukan banyak titik sampel
\item Banyak titik sampel dari pelemparan dua buah dadu secara bersamaan adalah \dots
    \begin{enumerate}[label=\alph*.]
        \item 6
        \item 12
        \item 24
        \item 36
        \item 64
    \end{enumerate}

\item Banyak titik sampel dari pelemparan 3 keping uang logam adalah \dots
    \begin{enumerate}[label=\alph*.]
        \item 3
        \item 6
        \item 8
        \item 9
        \item 12
    \end{enumerate}

% 23. Peluang
\item Peluang munculnya mata dadu ganjil pada pelemparan sebuah dadu adalah \dots
    \begin{enumerate}[label=\alph*.]
        \item $\frac{1}{6}$
        \item $\frac{1}{3}$
        \item $\frac{1}{2}$
        \item $\frac{2}{3}$
        \item $1$
    \end{enumerate}

\item Pada pelemparan sekeping uang logam, peluang muncul angka adalah \dots
    \begin{enumerate}[label=\alph*.]
        \item $\frac{1}{4}$
        \item $\frac{1}{3}$
        \item $\frac{1}{2}$
        \item $\frac{2}{3}$
        \item $1$
    \end{enumerate}

% 24. Frekuensi harapan
\item Jika sebuah dadu dilempar sebanyak 60 kali, frekuensi harapan munculnya mata dadu 4 adalah \dots
    \begin{enumerate}[label=\alph*.]
        \item 10 kali
        \item 15 kali
        \item 20 kali
        \item 30 kali
        \item 60 kali
    \end{enumerate}

\item Peluang seorang anak lulus ujian adalah 0,8. Jika ada 50 anak yang ikut ujian, frekuensi harapan anak yang lulus adalah \dots
    \begin{enumerate}[label=\alph*.]
        \item 10 anak
        \item 20 anak
        \item 30 anak
        \item 40 anak
        \item 50 anak
    \end{enumerate}

% 25. Kombinasi
\item Nilai dari kombinasi $C(5, 2)$ adalah \dots
    \begin{enumerate}[label=\alph*.]
        \item 5
        \item 10
        \item 15
        \item 20
        \item 60
    \end{enumerate}

\item Dari 7 siswa akan dipilih 3 siswa untuk mewakili lomba. Banyak cara pemilihannya adalah \dots
    \begin{enumerate}[label=\alph*.]
        \item 21
        \item 35
        \item 70
        \item 120
        \item 210
    \end{enumerate}

\item Perbedaan utama permutasi dan kombinasi adalah pada permutasi \dots, sedangkan pada kombinasi tidak.
    \begin{enumerate}[label=\alph*.]
        \item Memperhatikan warna
        \item Memperhatikan urutan
        \item Melibatkan dadu
        \item Menggunakan faktorial
        \item Melibatkan uang logam
    \end{enumerate}

\end{enumerate}
\end{document}
